\begin{figure}[ht]
\centering
\begin{tikzpicture}[>=Stealth, font=\small, node distance=1.35cm]
  \node[draw, rounded corners, fill=red!8, minimum width=5.8cm, minimum height=0.95cm] (flt)
    {假设 $a^p+b^p=c^p$ 有非平凡原始解，$p\ge 5$};
  \node[draw, rounded corners, fill=orange!12, below=of flt, minimum width=5.8cm, minimum height=1cm] (frey)
    {Frey 曲线 $E_{a,b,c}: y^2=x(x-a^p)(x+b^p)$};
  \node[draw, rounded corners, fill=yellow!15, below=of frey, minimum width=6.6cm, minimum height=1cm] (rep)
    {$\rhobar_{E,p}:G_\Q\to \GL_2(\F_p)$，半稳定且在 $q\mid abc$ 处可降级};
  \node[draw, rounded corners, fill=green!10, below=of rep, minimum width=6.6cm, minimum height=1cm] (wiles)
    {模性定理：$E_{a,b,c}$ 模化，故 $\rhobar_{E,p}$ 来自权 $2$ 新形式，级别 $2\,\mathrm{rad}(abc)$};
  \node[draw, rounded corners, fill=blue!8, below=of wiles, minimum width=6.6cm, minimum height=1cm] (ribet)
    {Ribet 水平下降：反复去掉每个 $q\mid abc$，降到级别 $2$};
  \node[draw, rounded corners, fill=purple!10, below=of ribet, minimum width=5.8cm, minimum height=0.95cm] (zero)
    {$S_2(\GammaO(2))=0$，不存在这样的新形式};
  \node[draw, rounded corners, fill=gray!12, below=of zero, minimum width=4.2cm, minimum height=0.95cm] (contra)
    {矛盾，故 FLT 成立};

  \draw[->, very thick] (flt) -- node[right] {Frey} (frey);
  \draw[->, very thick] (frey) -- node[right] {挠点与 Galois 作用} (rep);
  \draw[->, very thick] (rep) -- node[right] {Wiles--Taylor} (wiles);
  \draw[->, very thick] (wiles) -- node[right] {Ribet} (ribet);
  \draw[->, very thick] (ribet) -- node[right] {维数公式} (zero);
  \draw[->, very thick] (zero) -- (contra);
\end{tikzpicture}
\caption{费马大定理证明中的 Galois 表示路线图：从 Frey 曲线出发，经模性与水平下降，最终落到空的空间 $S_2(\GammaO(2))$。}
\label{fig:ch09-roadmap}
\end{figure}
